% ==============================================================================
% CAPÍTULO 1 - INTRODUÇÃO
% ==============================================================================

\chapter{Introdução}
\label{cap:introducao}

% Contextualização do tema

A escrita da introdução deve ser encarada como a construção de um roteiro lógico 
que guia o leitor desde o cenário geral até a sua contribuição específica. 

O texto deve iniciar obrigatoriamente pela delimitação do assunto, 
estabelecendo o contexto da pesquisa. Deve-se abordar, de forma sucinta, 
a grande área em que o trabalho está inserido e afunilar gradualmente para 
o tema específico, situando-o com clareza dentro do campo de estudo.

A partir desse cenário, é fundamental apresentar o problema central e os argumentos 
que justificam sua relevância, deixando nítido por que ele merece ser investigado. 

Na sequência, deve ser apresentada a solução proposta e os objetivos da pesquisa, 
demonstrando como essa abordagem avança o estado da arte e quais são 
os diferenciais que a tornam eficaz frente aos desafios existentes. 

Por fim, para consolidar a ambientação do leitor, a introdução deve apresentar 
de forma sucinta a estrutura do artigo, descrevendo o que será abordado em 
cada seção subsequente, garantindo assim uma transição fluida para o 
desenvolvimento do trabalho.
